% ======================================================================
% Chapter 10: Systematics of Renormalization
% (book pages 315-345; PDF pages 338-369)
% ======================================================================
\chapter{Systematics of Renormalization}

While computing radiative corrections in Chapters~6 and~7, we encountered
three QED diagrams with ultraviolet divergences: the electron self-energy, the
vertex correction, and the vacuum polarization.  In each case we saw that the
divergence could be regulated and canceled, yielding finite expressions for
measurable quantities.  In Chapter~8, we pointed out that such ultraviolet
divergences occur commonly and, in fact, naturally in quantum field theory
calculations.  We sketched a physical interpretation of these divergences, with
implications both in quantum field theory and in the statistical theory of
phase transitions.  In the next few chapters, we will convert this sketchy
picture into a quantitative theory that allows precise calculations.

In this chapter, we begin this study by developing a classification of the
ultraviolet divergences that can appear in a quantum field theory.  Rather
than stumbling across these divergences one by one and repairing them case by
case, we now set out to determine once and for all which diagrams are
divergent, and in which theories these divergences can be eliminated
systematically.  As examples we will consider both QED and scalar field
theories.

\section{Counting of Ultraviolet Divergences}

In this section we will use elementary arguments to determine, tentatively,
when a Feynman diagram contains an ultraviolet divergence.  We begin by
analyzing quantum electrodynamics.

First we introduce the following notation, to characterize a typical diagram
in QED:
\begin{equation}
  \begin{aligned}
    N_e &= \text{number of external electron lines};\\
    N_\gamma &= \text{number of external photon lines};\\
    P_e &= \text{number of electron propagators};\\
    P_\gamma &= \text{number of photon propagators};\\
    V &= \text{number of vertices};\\
    L &= \text{number of loops}.
  \end{aligned}
\end{equation}
(This analysis applies to correlation functions as well as scattering
amplitudes.  In the former case, propagators that are connected to external
points should be counted as external lines, not as propagators.)

The expression corresponding to a typical diagram looks like this:
\begin{equation}
  \int d^4k_1\,d^4k_2\cdots d^4k_L\
    \frac{(\slashed{k}_i\cdots)(\cdots)}{(k_i^2-m^2)\cdots(k_j^2)\cdots}
    \cdots .
\end{equation}
For each loop there is a potentially divergent 4-momentum integral, but each
propagator aids the convergence of this integral by putting one or two powers
of momentum into the denominator.  Very roughly speaking, the diagram
diverges unless there are more powers of momentum in the denominator than in
the numerator.  Let us therefore define the \emph{superficial degree of
divergence}, $D$, as the difference:
\begin{equation}
  D = (\text{power of }k\text{ in numerator}) - (\text{power of }k
    \text{ in denominator}) = 4L - P_e - 2P_\gamma .
  \label{eq:10.1}
\end{equation}
Naively, we expect a diagram to have a divergence proportional to
$\Lambda^D$, where $\Lambda$ is a momentum cutoff, when $D > 0$.  We expect a
divergence of the form $\log\Lambda$ when $D = 0$, and no divergence when
$D < 0$.

This naive expectation is often wrong, for one of three reasons (see
Fig.~\ref{fig:10.1}).  When a diagram contains a divergent subdiagram, its
actual divergence may be worse than that indicated by $D$.  When symmetries
(such as the Ward identity) cause certain terms to cancel, the divergence of a
diagram may be reduced or even eliminated.  Finally, a trivial diagram with no
propagators and no loops has $D = 0$ but no divergence.

\begin{figure}[t]
  \centering
  \includegraphics[width=0.6\textwidth]{fig10.1.png}
  \caption{Some simple QED diagrams that illustrate the superficial degree of
    divergence.  The first diagram is finite, even though $D = 0$.  The third
    diagram has $D = 2$ but only a logarithmic divergence, due to the Ward
    identity (see Section~7.5).  The fourth diagram diverges, even though
    $D < 0$, since it contains a divergent subdiagram.  Only in the second and
    fifth diagrams does the superficial degree of divergence coincide with the
    actual degree of divergence.}
  \label{fig:10.1}
\end{figure}

Despite all of these complications, $D$ is still a useful quantity.  To see
why, let us rewrite it in terms of the number of external lines
($N_e$, $N_\gamma$) and vertices ($V$).  Note that the number of loop
integrations in a diagram is
\begin{equation}
  L = P_e + P_\gamma - V + 1,
  \label{eq:10.2}
\end{equation}
since in our original Feynman rules each propagator has a momentum integral,
each vertex has a delta function, and one delta function merely enforces
overall momentum conservation.  Furthermore, the number of vertices is
\begin{equation}
  V = 2P_\gamma + N_\gamma = \tfrac{1}{2}(2P_e + N_e),
  \label{eq:10.3}
\end{equation}
since each vertex involves exactly one photon line and two electron lines.
(The propagators count twice since they have two ends on vertices.)  Putting
these relations together, we find that $D$ can be expressed as
\begin{equation}
  D = 4L - P_e - 2P_\gamma = 4 - \frac{3}{2}N_e - N_\gamma ,
  \label{eq:10.4}
\end{equation}
independent of the number of vertices.  The superficial degree of divergence
of a QED diagram depends only on the number of external legs of each type.

According to result \eqref{eq:10.4}, only diagrams with a small number of
external legs have $D \ge 0$; those seven types of diagrams are shown in
Fig.~\ref{fig:10.2}.  Since external legs do not enter the potentially
divergent integral, we can restrict our attention to amputated diagrams.  We
can also restrict our attention to one-particle-irreducible diagrams, since
reducible diagrams are simple products of the integrals corresponding to
their irreducible parts.  Thus the task of enumerating all of the divergent
QED diagrams reduces to that of analyzing the seven types of amputated,
one-particle-irreducible amplitudes shown in Fig.~\ref{fig:10.2}.  Other
diagrams may diverge, but only when they contain one of these seven as a
subdiagram.  Let us therefore consider each of these seven amplitudes in
turn.

\begin{figure}[t]
  \centering
  \includegraphics[width=0.6\textwidth]{fig10.2.png}
  \caption{The seven QED amplitudes whose superficial degree of divergence
    ($D$) is $\ge 0$.  (Each circle represents the sum of all possible QED
    diagrams.)  As explained in the text, amplitude (a) is irrelevant to
    scattering processes, while amplitudes (b) and (d) vanish because of
    symmetries.  Amplitude (e) is nonzero, but its divergent parts cancel due
    to the Ward identity.  The remaining amplitudes (c, f, and g) are all
    logarithmically divergent, even though $D > 0$ for (c) and (f).}
  \label{fig:10.2}
\end{figure}

The zero-point function, Fig.~\ref{fig:10.2}(a), is very badly divergent.  But
this object merely causes an unobservable shift of the vacuum energy; it never
contributes to $S$-matrix elements.

To analyze the photon one-point function (Fig.~\ref{fig:10.2}b), note that the
external photon must be attached to a QED vertex.  Neglecting the external
photon propagator, this amplitude is therefore
\begin{equation}
  \mathcal{M}^\mu(q) = -ie\int d^4x\ e^{-iq\cdot x}\
    \langle\Omega|T j^\mu(x)|\Omega\rangle ,
  \label{eq:10.5}
\end{equation}
where $j^\mu = \bar{\psi}\gamma^\mu\psi$ is the electromagnetic current
operator.  But the vacuum expectation value of $j^\mu$ must vanish by Lorentz
invariance, since otherwise it would be a preferred 4-vector.

The photon one-point function also vanishes for a second reason:
charge-conjugation invariance.  Recall that $C$ is a symmetry of QED, so
$C|\Omega\rangle = |\Omega\rangle$.  But $j^\mu(x)$ changes sign under charge
conjugation, $Cj^\mu(x)C^{\dagger} = -j^\mu(x)$, so its vacuum expectation
value must vanish:
\begin{equation}
  \langle\Omega|T j^\mu(x)|\Omega\rangle
  = \langle\Omega|C^{\dagger}Cj^\mu(x)C^{\dagger}C|\Omega\rangle
  = -\langle\Omega|T j^\mu(x)|\Omega\rangle = 0 .
\end{equation}
The same argument applies to any vacuum expectation value of an odd number of
electromagnetic currents.  In particular, the photon three-point function,
Fig.~\ref{fig:10.2}(d), vanishes.  (This result is known as \emph{Furry's
theorem}.)  It is not hard to check explicitly that the photon one- and
three-point functions vanish in the leading order of perturbation theory (see
Problem~10.1).

The remaining amplitudes in Fig.~\ref{fig:10.2} are all nonzero, so we must
analyze their structures in more detail.  Consider, for example, the electron
self-energy (Fig.~\ref{fig:10.2}f).  This amplitude is a function of the
electron momentum $p$, so let us expand it in a Taylor series about $p = 0$:
\begin{equation}
  \Sigma(\slashed{p}) = A_0 + A_1\slashed{p} + A_2\slashed{p}^2 + \cdots ,
\end{equation}
where each coefficient is independent of $p$:
\begin{equation}
  A_n = \frac{1}{n!}\ \frac{d^n}{d\slashed{p}^n}\ \Sigma(\slashed{p})
    \Big|_{\slashed{p}=0} .
\end{equation}
(These coefficients are infrared divergent; to compute them explicitly we
would need an infrared regulator, as in Chapter~6.)  The diagrams contributing
to the electron self-energy depend on $p$ through the denominators of
propagators.  To compute the coefficients $A_n$, we differentiate these
propagators, giving expressions like
\begin{equation}
  \frac{d}{dp_\mu}\ \frac{1}{\slashed{p}-m} = -\frac{1}{\slashed{p}-m}\,
    \gamma^\mu\ \frac{1}{\slashed{p}-m} .
\end{equation}
That is, each derivative with respect to the external momentum $p$ lowers the
superficial degree of divergence by 1.  Since the constant term $A_0$ has
(superficially) a linear divergence, $A_1$ can have only a logarithmic
divergence; all the remaining $A_n$ are finite.  (This argument breaks down
when the divergence is in a subdiagram, since then not all propagators involve
the large momentum $k$.  We will face this problem in Section~10.4.)

The electron self-energy amplitude has one additional subtlety.  If the
constant term $A_0$ were proportional to $\Lambda$ (the ultraviolet cutoff),
the electron mass shift would, according to the analysis in Section~7.1, also
have a term proportional to $\Lambda$.  But the electron mass shift must
actually be proportional to $m$, since chiral symmetry would forbid a mass
shift if $m_0$ were zero.  At worst, the constant term can be proportional to
$m\log\Lambda$.  We therefore expect the entire self-energy amplitude to have
the form
\begin{equation}
  -a_0\,m\log\Lambda + a_1\slashed{p}\log\Lambda + (\text{finite terms}),
  \label{eq:10.6}
\end{equation}
exactly what we found for the term of order $\alpha$ in Eq.~\eqref{eq:7.19}.

Let us analyze the exact electron-photon vertex, Fig.~\ref{fig:10.2}(g), in
the same way.  (Again we implicitly assume that infrared divergences have been
regulated.)  Expanding in powers of the three external momenta, we immediately
see that only the constant term is divergent, since differentiating with
respect to any external momentum would lower the degree of divergence to
$-1$.  This amplitude therefore contains only one divergent constant:
\begin{equation}
  \propto -ie\log\Lambda + \text{finite terms} .
  \label{eq:10.7}
\end{equation}
As discussed in Section~7.5, the photon self-energy (Fig.~\ref{fig:10.2}c) is
constrained by the Ward identity to have the form
\begin{equation}
  \Pi^{\mu\nu}(q) = (g^{\mu\nu}q^2 - q^\mu q^\nu)\,\Pi(q^2) .
  \label{eq:10.8}
\end{equation}
Viewing this expression as a Taylor series in $q$, we see that the constant and
linear terms both vanish, lowering the superficial degree of divergence from 2
to 0.  The only divergence, therefore, is in the constant term of
$\Pi(q^2)$, and this divergence is only logarithmic.  This result is exactly
what we found for the lowest-order contribution to $\Pi(q^2)$ in
Eq.~(7.90).

Finally, consider the photon-photon scattering amplitude,
Fig.~\ref{fig:10.2}(e).  The Ward identity requires that if we replace any
external photon by its momentum vector, the amplitude vanishes:
\begin{equation}
  k_{\mu_1}\mathcal{M}^{\mu_1\mu_2\mu_3\mu_4}(k_1,k_2,k_3,k_4) = 0 .
  \label{eq:10.9}
\end{equation}
By exhaustion one can show that this condition is satisfied only if the
amplitude is proportional to $(g^{\mu_1\mu_3}k^{\mu_2} - g^{\mu_1\mu_2}k^{\mu_3})$,
with a similar factor for each of the other three legs.  Each of these factors
involves one power of momentum, so all terms with less than four powers of
momentum in the Taylor series of this amplitude must vanish.  The first
nonvanishing term has $D = 0 - 4 = -4$, and therefore this amplitude is
finite.

In summary, we have found that there are only three ``primitively'' divergent
amplitudes in QED: the three that we already found in Chapters~6 and~7.
(Other amplitudes may also be divergent, but only because of diagrams that
contain these primitive amplitudes as components.)  Furthermore, the
dependence of these divergent amplitudes on external momenta is extremely
simple.  If we expand each amplitude as a power series in its external
momenta, there are altogether only four divergent coefficients in the
expansions.  In other words, QED contains only four divergent numbers.  In
the next section we will see how these numbers can be absorbed into
unobservable Lagrangian parameters, so that observable scattering amplitudes
are always finite.

For the remainder of this section, let us try to understand the superficial
degree of divergence from a more general viewpoint.  The theory of QED in four
spacetime dimensions is rather special, so let us first generalize to QED in
$d$ dimensions.  In this case, $D$ is given by
\begin{equation}
  D = dL - P_e - 2P_\gamma ,
  \label{eq:10.10}
\end{equation}
since each loop contributes a $d$-dimensional momentum integral.  Relations
\eqref{eq:10.2} and \eqref{eq:10.3} still hold, so we can again rewrite $D$ in
terms of $V$, $N_e$, and $N_\gamma$:
\begin{equation}
  D = d - \frac{d-1}{2}N_e - \frac{d-2}{2}N_\gamma .
\end{equation}
For $d = 4$, this reduces to \eqref{eq:10.4}.  For $d > 4$, the divergences
become worse; for $d < 4$, they are milder.

Let us now perform the analogous analysis for a theory of a real scalar field
with a $\phi^n$ interaction in $d$ spacetime dimensions.  The Lagrangian is
\begin{equation}
  \mathcal{L} = \tfrac{1}{2}(\partial_\mu\phi)^2 - \tfrac{1}{2}m^2\phi^2
    - \frac{\lambda}{n!}\phi^n .
  \label{eq:10.12}
\end{equation}
The superficial degree of divergence of a diagram with $N$ external lines is
\begin{equation}
  D = dL - 2P ,
  \label{eq:10.13}
\end{equation}
where $P$ is the number of scalar propagators.  Using
$L = P - V + 1$ and $N = nV - 2P$, we can rewrite $D$ as
\begin{equation}
  D = d + V\Bigl[ \frac{n(d-2)}{2} - d \Bigr] - \frac{d-2}{2}N .
  \label{eq:10.13b}
\end{equation}
The theory is renormalizable if the coefficient of $V$ is zero, that is, if
\begin{equation}
  n = \frac{2d}{d-2} .
\end{equation}
For $d = 4$, this gives $n = 4$: In four dimensions a $\phi^4$ coupling is
renormalizable, while higher powers of $\phi$ are non-renormalizable.  In three
dimensions a $\phi^6$ coupling becomes renormalizable, while $\phi^4$ is
super-renormalizable.  In two spacetime dimensions any coupling of the form
$\phi^n$ is super-renormalizable.

Expression \eqref{eq:10.13} can also be derived in a somewhat different way,
from dimensional analysis.  In any quantum field theory, the action
$S = \int d^dx\ \mathcal{L}$ must be dimensionless, since we work in units
where $\hbar = 1$.  In this system of units, the integral $d^dx$ has units
$(\mathrm{mass})^{-d}$, and so the Lagrangian has units
$(\mathrm{mass})^d$.  Since all units can be expressed as powers of mass, it is
unambiguous to say simply that the Lagrangian has ``dimension $d$.''  Using
this result, we can infer from the explicit form of \eqref{eq:10.12} the
dimensions of the field $\phi$ and the coupling constant $\lambda$.  From the
kinetic term in $\mathcal{L}$ we see that $\phi$ has dimension $(d-2)/2$.  Note
that the parameter $m$ consistently has dimensions of mass.  From the
interaction term and the dimension of $\phi$, we infer that the $\lambda$ has
dimension
\begin{equation}
  [\lambda] = d - n\frac{d-2}{2} .
\end{equation}
Now consider an arbitrary diagram with $N$ external lines.  One way that such
a diagram could arise is from an interaction term $\eta\phi^N$ in the
Lagrangian.  The dimension of $\eta$ would then be
$d - N(d-2)/2$, and therefore we conclude that any (amputated) diagram with
$N$ external lines has dimension $d - N(d-2)/2$.  In our theory with only the
$\lambda\phi^n$ vertex, if the diagram has $V$ vertices, its divergent part is
proportional to $\lambda^V\Lambda^D$, where $\Lambda$ is a high-momentum
cutoff and $D$ is the superficial degree of divergence.  (This is the
``generic'' case; all the exceptions noted above also apply here.)  Applying
dimensional analysis, we find
\begin{equation}
  [\lambda^V\Lambda^D] = V\Bigl[ d - \frac{n(d-2)}{2} \Bigr] + D
  = d - \frac{N(d-2)}{2},
\end{equation}
in agreement with \eqref{eq:10.13}.

Note that the quantity that multiplies $V$ in this expression is just the
dimension of the coupling constant $\lambda$.  This analysis can be carried out
for QED and other field theories, with the same result.  Thus we can
characterize the three degrees of renormalizability in a second way:
\begin{equation}
  \begin{aligned}
    \text{Super-Renormalizable:}\quad & \text{Coupling constant has positive
      mass dimension.}\\
    \text{Renormalizable:}\quad & \text{Coupling constant is dimensionless.}\\
    \text{Non-Renormalizable:}\quad & \text{Coupling constant has negative
      mass dimension.}
  \end{aligned}
\end{equation}
This is exactly the conclusion that we stated without proof in Section~4.1.
In QED, the coupling constant $e$ is dimensionless; thus QED is (at least
superficially) renormalizable.

\section{Renormalized Perturbation Theory}

In the previous section we saw that a renormalizable quantum field theory
contains only a small number of superficially divergent amplitudes.  In QED,
for example, there are three such amplitudes, containing four infinite
constants.  In Chapters~6 and~7 these infinities disappeared by the end of our
computations: The infinity in the vertex correction diagram was canceled by
the electron field-strength renormalization, while the infinity in the vacuum
polarization diagram caused only an unobservable shift of the electron's
charge.  In fact, it is generally true that the divergences in a renormalizable
quantum field theory never show up in observable quantities.

To obtain a finite result for an amplitude involving divergent diagrams, we
have so far used the following procedure: Compute the diagrams using a
regulator, to obtain an expression that depends on the bare mass ($m_0$), the
bare coupling constant ($e_0$), and some ultraviolet cutoff ($\Lambda$).
Then compute the physical mass ($m$) and the physical coupling constant
($e$), to whatever order is consistent with the rest of the calculation; these
quantities will also depend on $m_0$, $e_0$, and $\Lambda$.  To calculate an
$S$-matrix element (rather than a correlation function), one must also compute
the field-strength renormalization(s) $Z$ (in accord with Eq.~\eqref{eq:7.45}).
Combining all of these expressions, eliminate $m_0$ and $e_0$ in favor of $m$
and $e$; this step is the ``renormalization.''  The resulting expression for
the amplitude should be finite in the limit $\Lambda\to\infty$.

The above procedure always works in a renormalizable quantum field theory.
However, it can often be cumbersome, especially at higher orders in
perturbation theory.  In this section we will develop an alternative procedure
which works more automatically.  We will do this first for $\phi^4$ theory,
returning to QED in the next section.

The Lagrangian of $\phi^4$ theory is
\begin{equation}
  \mathcal{L} = \tfrac{1}{2}(\partial_\mu\phi)^2 - \tfrac{1}{2}m_0^2\phi^2
    - \frac{\lambda_0}{4!}\phi^4 .
\end{equation}
We now write $m_0$ and $\lambda_0$, to emphasize that these are the bare values
of the mass and coupling constant, not the values measured in experiments.

The superficial degree of divergence of a diagram with $N$ external legs is,
according to \eqref{eq:10.13},
\begin{equation}
  D = 4 - N .
\end{equation}
Since the theory is invariant under $\phi\to-\phi$, all amplitudes with an odd
number of external legs vanish.  The only divergent amplitudes are therefore
the two-point function,
\begin{equation}
  \sim \Lambda^2 + p^2\log\Lambda + (\text{finite terms}) ;
\end{equation}
the four-point function,
\begin{equation}
  \sim \log\Lambda + (\text{finite terms}) ;
\end{equation}
and the vacuum energy,
\begin{equation}
  \text{(unobservable vacuum energy shift)} .
\end{equation}
Ignoring the vacuum diagram, these amplitudes contain three infinite
constants.  Our goal is to absorb these constants into the three unobservable
parameters of the theory: the bare mass, the bare coupling constant, and the
field strength.  To accomplish this goal, it is convenient to reformulate the
perturbation expansion so that these unobservable quantities do not appear
explicitly in the Feynman rules.

First we will eliminate the shift in the field strength.  Recall from
Section~7.1 that the exact two-point function has the form
\begin{equation}
  \int d^4x\ \langle\Omega|T\phi(x)\phi(0)|\Omega\rangle\ e^{ip\cdot x}
  = \frac{iZ}{p^2-m^2} + (\text{terms regular at } p^2 = m^2),
  \label{eq:10.14}
\end{equation}
where $m$ is the physical mass.  We can eliminate the awkward residue $Z$ from
this equation by rescaling the field:
\begin{equation}
  \phi = Z^{1/2}\phi_r .
  \label{eq:10.15}
\end{equation}
This transformation changes the values of correlation functions by a factor of
$Z^{-1/2}$ for each field.  Thus, in computing $S$-matrix elements, we no
longer need the factors of $Z$ in Eq.~\eqref{eq:7.45}; a scattering amplitude
is simply the sum of all connected, amputated diagrams, exactly as we
originally guessed in Eq.~\eqref{eq:4.103}.

The Lagrangian is much uglier after the rescaling:
\begin{equation}
  \mathcal{L} = \tfrac{1}{2}Z(\partial_\mu\phi_r)^2
    - \tfrac{1}{2}m_0^2Z\phi_r^2 - \frac{\lambda_0}{4!}Z^2\phi_r^4 .
  \label{eq:10.16}
\end{equation}
The bare mass and coupling constant still appear in $\mathcal{L}$, but they can
be eliminated as follows.  Define
\begin{equation}
  \delta_Z = Z - 1,\qquad
  \delta_m = m_0^2Z - m^2,\qquad
  \delta_\lambda = \lambda_0Z^2 - \lambda,
  \label{eq:10.17}
\end{equation}
where $m$ and $\lambda$ are the physically measured mass and coupling
constant.  Then the Lagrangian becomes
\begin{equation}
  \mathcal{L} = \tfrac{1}{2}(\partial_\mu\phi_r)^2
    - \tfrac{1}{2}m^2\phi_r^2 - \frac{\lambda}{4!}\phi_r^4
    + \tfrac{1}{2}\delta_Z(\partial_\mu\phi_r)^2
      - \tfrac{1}{2}\delta_m\phi_r^2 - \frac{\delta_\lambda}{4!}\phi_r^4 .
  \label{eq:10.18}
\end{equation}
The first line now looks like the familiar $\phi^4$-theory Lagrangian, but is
written in terms of the physical mass and coupling.  The terms in the second
line, known as \emph{counterterms}, have absorbed the infinite but
unobservable shifts between the bare parameters and the physical parameters.
It is tempting to say that we have ``added'' these counterterms to the
Lagrangian, but in fact we have merely split each term in \eqref{eq:10.16}
into two pieces.

The definitions in \eqref{eq:10.17} are not useful unless we give precise
definitions of the physical mass and coupling constant.  Equation
\eqref{eq:10.14} defines $m^2$ as the location of the pole in the propagator.
There is no obviously best definition of $\lambda$, but a perfectly good
definition would be obtained by setting $\lambda$ equal to the magnitude of the
scattering amplitude at zero momentum.  Thus we have the two defining
relations,
\begin{equation}
  \frac{i}{p^2-m^2} + (\text{terms regular at } p^2 = m^2) ;
  \qquad -i\lambda = \text{(amputated amplitude at } s = 4m^2,\
    t = u = 0) .
  \label{eq:10.19}
\end{equation}
These equations are called \emph{renormalization conditions}.  (The first
equation actually contains two conditions, specifying both the location of the
pole and its residue.)

Our new Lagrangian, Eq.~\eqref{eq:10.18}, gives a new set of Feynman rules,
shown in Fig.~\ref{fig:10.3}.  The propagator and the first vertex come from
the first line of \eqref{eq:10.18}, and are identical to the old rules except
for the appearance of the physical mass and coupling in place of the bare
values.  The counterterms in the second line of \eqref{eq:10.18} give two new
vertices (also called counterterms):
\begin{equation}
  \text{(propagator counterterm)} = i(p^2\delta_Z - \delta_m);
  \qquad \text{(four-point counterterm)} = -i\delta_\lambda .
  \label{eq:10.20}
\end{equation}

\begin{figure}[t]
  \centering
  \includegraphics[width=0.6\textwidth]{fig10.3.png}
  \caption{Feynman rules for $\phi^4$ theory in renormalized perturbation
    theory.}
  \label{fig:10.3}
\end{figure}

We can use these new Feynman rules to compute any amplitude in $\phi^4$
theory.  The procedure is as follows.  Compute the desired amplitude as the
sum of all possible diagrams created from the propagator and vertices shown in
Fig.~\ref{fig:10.3}.  The loop integrals in the diagrams will often diverge, so
one must introduce a regulator.  The result of this computation will be a
function of the three unknown parameters $\delta_Z$, $\delta_m$, and
$\delta_\lambda$.  Adjust (or ``renormalize'') these three parameters as
necessary to maintain the renormalization conditions \eqref{eq:10.19}.  After
this adjustment, the expression for the amplitude should be finite and
independent of the regulator.

This procedure, using Feynman rules with counterterms, is known as
\emph{renormalized perturbation theory}.  It should be contrasted with the
procedure we used in Part~I, outlined at the beginning of this section, which
is called \emph{bare perturbation theory} (since the Feynman rules involve the
bare mass and coupling constant).  The two methods are completely equivalent.
The differences between them are purely a matter of bookkeeping.  You will get
the same answers using either procedure, so you may choose whichever you find
more convenient.  In general, renormalized perturbation theory is technically
easier to use, especially for multiloop diagrams; however, bare perturbation
theory is sometimes easier for complicated one-loop calculations.  We will use
renormalized perturbation theory in most of the rest of this book.

\subsection{One-Loop Structure of $\phi^4$ Theory}

To make more sense of the renormalization procedure, let us carry it out
explicitly at the one-loop level.

First consider the basic two-particle scattering amplitude,
\begin{equation}
  i\mathcal{M}(p_1p_2\to p_3p_4) = \text{(tree)} + \text{(one-loop)}
    + \cdots .
\end{equation}
If we define $p = p_1+p_2$, then the second diagram is
\begin{equation}
  = (-i\lambda)^2\cdot iV(p^2) ,
  \label{eq:10.21}
\end{equation}
where
\begin{equation}
  V(p^2) = \frac{1}{2}\int \frac{d^dk}{(2\pi)^d}\ \frac{1}{k^2-m^2+i\epsilon}\
    \frac{1}{(k+p)^2-m^2+i\epsilon} .
  \label{eq:10.22}
\end{equation}
Note that $p^2$ is equal to the Mandelstam variable $s$.  The next two diagrams
are identical, except that $s$ will be replaced by $t$ and $u$.  The entire
amplitude is therefore
\begin{equation}
  i\mathcal{M} = -i\lambda + (-i\lambda)^2\bigl[ iV(s) + iV(t) + iV(u) \bigr]
    - i\delta_\lambda .
  \label{eq:10.23}
\end{equation}
According to our renormalization condition \eqref{eq:10.19}, this amplitude
should be equal to $-i\lambda$ when $s = 4m^2$, $t = u = 0$.  This fixes
$\delta_\lambda$:
\begin{equation}
  \delta_\lambda = (-\lambda)^2\bigl[ V(4m^2) + 2V(0) \bigr] .
  \label{eq:10.24}
\end{equation}
Substituting this back into \eqref{eq:10.23}, we obtain the finite result
\begin{equation}
  i\mathcal{M} = -i\lambda + (-i\lambda)^2\bigl[ iV(s) - iV(4m^2)
    + iV(t) - iV(0) + iV(u) - iV(0) \bigr] ,
\end{equation}
where the differences $V(s) - V(4m^2)$ and $V(t) - V(0)$ are finite, since the
divergences are independent of the external momenta.

Next consider the two-point function at one loop.  The one-loop correction to
the propagator is
\begin{equation}
  = (-i\lambda)\ \frac{1}{2}\int \frac{d^dk}{(2\pi)^d}\ \frac{1}{k^2-m^2}
  = -i\lambda\,V_0 ,
\end{equation}
where $V_0$ is the value of the loop integral with both external momenta zero.
This diagram is quadratically divergent.  Together with the counterterm, the
exact two-point function is
\begin{equation}
  = \frac{i}{p^2-m^2} + \frac{i}{p^2-m^2}\ \bigl[
    i(p^2\delta_Z - \delta_m) - i\lambda V_0 \bigr]\ \frac{i}{p^2-m^2}
    + \cdots .
  \label{eq:10.28}
\end{equation}
The renormalization conditions \eqref{eq:10.19} require that the pole remain
at $p^2 = m^2$ with residue $i$.  This fixes the counterterms
\begin{equation}
  \delta_m = -\lambda V_0,\qquad \delta_Z = 0 .
  \label{eq:10.29}
\end{equation}
Thus, in $\phi^4$ theory at one loop, the field-strength renormalization
counterterm $\delta_Z$ vanishes.  (This is a special feature of $\phi^4$
theory; the vanishing of $\delta_Z$ at one-loop order does not occur in more
general theories of scalar fields.)

The Yukawa theory described in Section~4.7 gives an explicit example of a
one-loop correction for which this counterterm is required.  In the Yukawa
theory, the scalar field propagator receives corrections at order $g^2$ from a
fermion loop diagram and the two propagator counterterms.  Using the Feynman
rules to compute the loop diagram, we find
\begin{equation}
  \begin{aligned}
    -iM^2(p^2) &= (-ig)^2\ \int \frac{d^dk}{(2\pi)^d}\ \mathrm{tr}\Bigl[
      \frac{i(\slashed{k}+m_f)}{k^2-m_f^2}\ \frac{i(\slashed{k}+\slashed{p}
        +m_f)}{(k+p)^2-m_f^2} \Bigr]\\
    &\quad + i(p^2\delta_Z - \delta_m) ,
  \end{aligned}
  \label{eq:10.32}
\end{equation}
where $m_f$ is the mass of the fermion that couples to the Yukawa field.  To
evaluate the integral, combine denominators and shift as in Eq.~\eqref{eq:10.23}.
Then the first term in the last line becomes
\begin{equation}
  \begin{aligned}
    &= 4g^2\int_0^1 dx\ \int \frac{d^d\ell}{(2\pi)^d}\
      \frac{\ell^2 - x(1-x)p^2 + m_f^2}{(\ell^2 + x(1-x)p^2 - m_f^2)^2}\\
    &= -4g^2\int_0^1 dx\ \frac{i}{(4\pi)^{d/2}}\ \frac{\Gamma(1-d/2)}
      {\Delta^{1-d/2}} + 4g^2\int_0^1 dx\ \frac{i(d-1)}{(4\pi)^{d/2}}\
      \frac{\Gamma(2-d/2)}{\Delta^{2-d/2}} ,
  \end{aligned}
  \label{eq:10.33}
\end{equation}
where $\Delta = m_f^2 - x(1-x)p^2$.

Now we can see that both of the counterterms $\delta_m$ and $\delta_Z$ must
take nonzero values in order to satisfy the renormalization conditions
\eqref{eq:10.28}.  To determine $\delta_m$, we subtract the value of the loop
diagram at $p^2 = m^2$ as before, so that
\begin{equation}
  \delta_m = -4g^2\int_0^1 dx\ \frac{i}{(4\pi)^{d/2}}\ \frac{\Gamma(2-d/2)}
    {\bigl[ m_f^2 - x(1-x)m^2 \bigr]^{2-d/2}} .
  \label{eq:10.34}
\end{equation}
To determine $\delta_Z$, we cancel also the first derivative with respect to
$p^2$ of the loop integral \eqref{eq:10.33}.  This gives
\begin{equation}
  \delta_Z = -\frac{4g^2(d-1)}{(4\pi)^{d/2}}\ \int_0^1 dx\ x(1-x)\
    \frac{\Gamma(2-d/2)}{\bigl[ m_f^2 - x(1-x)m^2 \bigr]^{2-d/2}} .
  \label{eq:10.35}
\end{equation}
Thus, in Yukawa theory, the propagator corrections at one-loop order require a
quadratically divergent mass renormalization and a logarithmically divergent
field strength renormalization.  This is the usual situation in scalar field
theories.

\section{Renormalization of Quantum Electrodynamics}

The procedure we followed in the previous section, yielding a ``renormalized''
perturbation theory formulated in terms of physically measurable parameters,
can be summarized as follows:
\begin{enumerate}
  \item Absorb the field-strength renormalizations into the Lagrangian by
    rescaling the fields.
  \item Split each term of the Lagrangian into two pieces, absorbing the
    infinite and unobservable shifts into counterterms.
  \item Specify the renormalization conditions, which define the physical
    masses and coupling constants and keep the field-strength renormalizations
    equal to 1.
  \item Compute amplitudes with the new Feynman rules, adjusting the
    counterterms as necessary to maintain the renormalization conditions.
\end{enumerate}
Let us now use this procedure to construct a renormalized perturbation theory
for Quantum Electrodynamics.

The original QED Lagrangian is
\begin{equation}
  \mathcal{L} = -\frac{1}{4}(F_{\mu\nu})^2
    + \bar{\psi}(i\slashed{\partial} - m_0)\psi - e_0\bar{\psi}\gamma^\mu\psi
      A_\mu .
\end{equation}
Computing the electron and photon propagators with this Lagrangian, we would
find expressions of the general form
\begin{equation}
  \frac{iZ_2}{\slashed{p}-m},\qquad \frac{-ig_{\mu\nu}Z_3}{q^2} .
\end{equation}
(We found just such expressions in the explicit one-loop calculations of
Chapter~7.)  To absorb $Z_2$ and $Z_3$ into $\mathcal{L}$, and hence eliminate
them from formula \eqref{eq:7.45} for the $S$-matrix, we substitute
$\psi = Z_2^{1/2}\psi_r$ and $A^\mu = Z_3^{1/2}A_r^\mu$.  Then the Lagrangian
becomes
\begin{equation}
  \mathcal{L} = -\frac{1}{4}Z_3(F_{\mu\nu}^r)^2 + Z_2\bar{\psi}_r
    (i\slashed{\partial} - m_0)\psi_r - e_0Z_2Z_3^{1/2}\bar{\psi}_r\gamma^\mu
      \psi_r A_r^\mu .
  \label{eq:10.36}
\end{equation}
We can introduce the physical electric charge $e$, measured at large distances
($q = 0$), by defining a scaling factor $Z_1$ as follows:%
\footnote{Since we define $e$ by the renormalization condition
$\Gamma^\mu(q=0) = \gamma^\mu$, the factor of $Z_1$ in the Lagrangian must
cancel the multiplicative correction factor that arises from loop corrections.
Therefore this definition of $Z_1$ is equivalent to that given in
Eq.~(7.47).}
\begin{equation}
  e_0Z_2Z_3^{1/2} = eZ_1 .
  \label{eq:10.37}
\end{equation}
If we let $m$ be the physical mass (the location of the pole in the electron
propagator), then we can split each term of the Lagrangian into two pieces as
follows:
\begin{equation}
  \begin{aligned}
    \mathcal{L} &= -\frac{1}{4}(F_{\mu\nu}^r)^2 + \bar{\psi}_r(i\slashed{\partial}
      - m)\psi_r - e\bar{\psi}_r\gamma^\mu\psi_r A_r^\mu\\
    &\quad - \frac{1}{4}\delta_3(F_{\mu\nu}^r)^2 + \bar{\psi}_r(i\delta_2\slashed{\partial}
      - \delta_m)\psi_r - e\delta_1\bar{\psi}_r\gamma^\mu\psi_r A_r^\mu ,
  \end{aligned}
  \label{eq:10.38}
\end{equation}
where
\begin{equation}
  \delta_3 = Z_3 - 1,\qquad
  \delta_m = Z_2m_0 - m,\qquad
  \delta_2 = Z_2 - 1,\qquad
  \delta_1 = Z_1 - 1 = (e_0/e)Z_2Z_3^{1/2} - 1 .
\end{equation}
The Feynman rules for renormalized QED are shown in Fig.~\ref{fig:10.4}.  In
addition to the familiar propagators and vertex, there are three counterterm
vertices.  The $\bar{e}e$ and $\bar{e}e\gamma$ counterterm vertices can be read
directly from the Lagrangian \eqref{eq:10.38}.  To derive the two-photon
counterterm, integrate $-\tfrac{1}{4}\delta_3(F_{\mu\nu})^2$ by parts to
obtain $\tfrac{1}{2}A^\mu\delta_3(g_{\mu\nu}\partial^2 - \partial_\mu
\partial_\nu)A^\nu$; this gives the expression shown in the figure.  In the
remainder of the book, when we set up renormalized perturbation theory, we
will drop the subscript $r$ used here to distinguish the rescaled fields.

\begin{figure}[t]
  \centering
  \includegraphics[width=0.6\textwidth]{fig10.4.png}
  \caption{Feynman rules for Quantum Electrodynamics in renormalized
    perturbation theory.}
  \label{fig:10.4}
\end{figure}

Each of the four counterterm coefficients must be fixed by a renormalization
condition.  The four conditions that we require have already been stated
implicitly: Two of them fix the electron and photon field-strength
renormalizations to 1, while the other two define the physical electron mass
and charge.  To write these conditions more explicitly, recall our notation
from Chapters~6 and~7:
\begin{equation}
  \begin{aligned}
    \frac{}{} &= i\Pi^{\mu\nu}(q) = i(g^{\mu\nu}q^2 - q^\mu q^\nu)\Pi(q^2),
    \qquad \Sigma(\slashed{p}) = i(\slashed{p}\delta_2 - \delta_m)
      + \cdots ,\\
    -ie\Gamma^\mu(p',p) &= \text{(amputated vertex)},
  \end{aligned}
  \label{eq:10.39}
\end{equation}
These amplitudes are now to be computed in renormalized perturbation theory;
that is, we are now redefining $\Pi(q^2)$, $\Sigma(\slashed{p})$, and
$\Gamma^\mu(p',p)$ to include counterterm vertices.  Furthermore, the new
definition of $\Gamma$ involves the physical electron charge.  With this
notation, the four conditions are
\begin{equation}
  \Sigma(\slashed{p}=m) = 0;\qquad
  \Pi(q^2=0) = 0;\qquad
  \frac{d\Sigma}{d\slashed{p}}\Big|_{\slashed{p}=m} = 0;\qquad
  -ie\Gamma^\mu(p'-p = 0) = -ie\gamma^\mu .
  \label{eq:10.40}
\end{equation}
The first condition fixes the electron mass at $m$, while the next two fix the
residues of the electron and photon propagators at 1.  Given these conditions,
the final condition fixes the electron charge to be $e$.

\subsection{One-Loop Structure of QED}

The four conditions \eqref{eq:10.40} allow us to determine the four
counterterms in \eqref{eq:10.38} in terms of the values of loop diagrams.  In
Chapters~6 and~7 we computed all of the diagrams required to carry out this
determination to one-loop order.  We will now collect these results and find
explicit expressions for the renormalization constants of QED to order
$\alpha$.  For overall consistency, we will use dimensional regularization,
which preserves the Ward identity, throughout.  The results are:
\begin{equation}
  \delta_2 = \frac{-\alpha}{2\pi}\int_0^1 dx\ (1-x)\ \Bigl[
    \frac{2}{\epsilon} - \gamma_E + \log 4\pi - \log(m^2-x(1-x)m^2) \Bigr]
    + \cdots ,
  \label{eq:10.41}
\end{equation}
\begin{equation}
  \delta_m = m\ \frac{-\alpha}{2\pi}\int_0^1 dx\ 2\ \Bigl[
    \frac{2}{\epsilon} - \gamma_E + \log 4\pi - \log(m^2-x(1-x)m^2) \Bigr]
    + \cdots ,
  \label{eq:10.42}
\end{equation}
\begin{equation}
  \delta_1 = \delta_2 \qquad \text{(by the Ward identity)},
  \label{eq:10.43}
\end{equation}
\begin{equation}
  \delta_3 = \frac{-\alpha}{3\pi}\ \Bigl[ \frac{2}{\epsilon} - \gamma_E
    + \log 4\pi - \log m^2 \Bigr] + \cdots .
  \label{eq:10.44}
\end{equation}
The equality $\delta_1 = \delta_2$, which follows from the Ward identity, is
the statement that the electron charge renormalization is entirely determined
by the field-strength renormalizations.  It is this relation that ensures that
the theory contains only three independent renormalization constants.

\section{Renormalization Beyond the Leading Order}

In Section~10.1 we counted the ultraviolet divergences of QED and $\phi^4$
theory and found only a small number of superficially divergent amplitudes.
The analysis of that section relied on a single-loop understanding of
divergences.  We must now face the problem that arises when we go to higher
orders in perturbation theory: diagrams can contain divergent subdiagrams,
and the naive counting based on the superficial degree of divergence of the
whole diagram can fail.

Let us consider the two-loop contribution to the four-point function in
$\phi^4$ theory.  The two-loop diagrams can be classified into those that
contain divergent subdiagrams (such as one-loop self-energy or vertex
corrections) and those that do not.  For the diagrams containing divergent
subdiagrams, the naive counting gives a divergence, but the subdivergence is
localized in a subdiagram that can be contracted to a point.  The philosophy of
renormalized perturbation theory is that these subdivergences should be
removed by the counterterms associated with the divergent subdiagrams.

A more difficult situation occurs when we have \emph{nested} or
\emph{overlapping} divergences, that is, when two divergent loops share a
propagator.  Some examples of diagrams with overlapping divergences are
\begin{equation}
  \text{(in } \phi^4 \text{ theory)}:\quad \text{(sunset diagram)};\qquad
  \text{(in QED)}:\quad \text{(two-loop vacuum polarization)} .
\end{equation}
To see the difficulty, consider the photon self-energy diagram.  One
contribution to this diagram comes from the region of momentum space where
$k_2$ is very large.  This means that, in position space, $x$, $y$, and $z$ are
very close together, while $w$ can be farther away.  In this region we can
think of the virtual photon as giving a correction to the vertex at $x$.  We
saw in Section~6.3 that this vertex correction is logarithmically divergent,
of the form
\begin{equation}
  \sim a\log\Lambda ,
\end{equation}
in the limit $\Lambda\to\infty$.  Plugging this vertex into the rest of the
diagram and integrating over $k_1$, we obtain an expression identical to the
one-loop photon self-energy correction $\Pi_2(q^2)$, displayed in
Eq.~(7.90), multiplied by the additional logarithmic divergence:
\begin{equation}
  \sim \alpha(g^{\mu\nu}q^2 - q^\mu q^\nu)\bigl( \log\Lambda^2
    + \log q^2 \bigr)\cdot\alpha\log\Lambda^2 .
  \label{eq:10.50}
\end{equation}
The $\log^2\Lambda^2$ term comes from the region where both $k_1$ and $k_2$
are large, while the $\log q^2\log\Lambda^2$ term comes from the region where
$k_2$ is large but $k_1$ is small.  Another such term would come from the
region where $k_1$ is large but $k_2$ is small.

The appearance of terms proportional to $\Pi_2(q^2)\cdot\log\Lambda^2$ in the
two-loop vacuum polarization diagram contradicts our naive argument, based on
the criterion of the superficial degree of divergence, that the divergent
terms of a Feynman integral are always simple polynomials in $q^2$.  We will
refer to divergences multiplying only polynomials in $q^2$ as \emph{local
divergences}, since their Fourier transforms back to position space are delta
functions or derivatives of delta functions.  We will call the new,
nonpolynomial, term a \emph{nonlocal divergence}.  Fortunately, our derivation
of the nonlocal divergent term gave this term a physical interpretation: It is
a local divergence surrounded by an ordinary, nondivergent, quantum field
theory process.

If this picture accurately describes all of the divergent terms of the
two-loop diagram, we should expect that these divergences are canceled by two
types of counterterm diagrams.  First, we can build diagrams of order
$\alpha^2$ by inserting the order-$\alpha$ counterterm vertex into the
one-loop vacuum polarization diagram.  These diagrams should cancel the
nonlocal divergence in \eqref{eq:10.50} and the corresponding contribution
from the region where $k_1$ is large and $k_2$ is small.  In fact, a detailed
analysis shows that the sum of the original diagram and these two counterterm
diagrams contains only local divergences.  Once these diagrams are added, the
only divergence that remains is a local one, which can be canceled by the
diagram that adds an order-$\alpha^2$ term to $\delta_3$.

We can extend the lessons of this example to a general picture of the
divergences of higher-loop Feynman diagrams and their cancellation.  A given
diagram may contain local divergences, as predicted by the analysis of
Section~10.1.  It may also contain nonlocal divergences due to divergent
subgraphs embedded in loops carrying small momenta.  These divergences are
canceled by diagrams in which the divergent subgraphs are replaced by their
counterterm vertices.  One might still ask two questions: First, does this
procedure remove all nonlocal divergences?  Second, does this procedure
preserve the finiteness of amplitudes, such as Eq.~(10.49), that are not
expected to be divergent by the superficial criteria of Section~10.1?  To
answer these questions requires an intricate study of nested Feynman
integrals.  The general analysis was begun by Bogoliubov and Parasiuk,
completed by Hepp, and elegantly refined by Zimmermann;%
\footnote{N.~N.~Bogoliubov and O.~S.~Parasiuk, \emph{Acta Math.} \textbf{97},
227 (1957); K.~Hepp, \emph{Comm.\ Math.\ Phys.} \textbf{2}, 301 (1966);
W.~Zimmermann, in Deser, et.\ al.\ (1970).} they showed that the answer to
both questions is yes.  Their result, known as the \emph{BPHZ theorem}, states
that, for a general renormalizable quantum field theory, to any order in
perturbation theory, all divergences are removed by the counterterm vertices
corresponding to superficially divergent amplitudes.  In other words, any
superficially renormalizable quantum field theory is in fact rendered finite
when one performs renormalized perturbation theory with the complete set of
counterterms.

The proof of the BPHZ theorem is quite technical, and we will not include it in
this book.  Instead, we will investigate one detailed example of a two-loop
calculation, which demonstrates explicitly the appearance and cancellation of
nonlocal divergences.

\section{A Two-Loop Example}

To illustrate the issues discussed in the previous section, let us consider the
two-loop contribution to the four-point function in $\phi^4$ theory.  There
are 16 relevant diagrams, shown in Fig.~\ref{fig:10.5}.  (There are also
several diagrams involving the one-loop correction to the propagator.  But
each of these is exactly canceled by its counterterm, as we saw in
Eq.~\eqref{eq:10.29}, so we can just ignore them.)  Fortunately, many of the
diagrams are simply related to each other.  Crossing symmetry reduces the
number of distinct diagrams to only six:
\begin{equation}
  \text{(six distinct two-loop diagrams)} .
  \label{eq:10.51}
\end{equation}
If this sum of diagrams is finite, then simply replacing $s$ with $t$ or $u$
gives a finite result for the remaining diagrams.

\begin{figure}[t]
  \centering
  \includegraphics[width=0.65\textwidth]{fig10.5.png}
  \caption{The two-loop contributions to the four-point function in $\phi^4$
    theory.  Note that the diagrams in the first three lines are related to
    each other by crossing, being in the $s$-, $t$-, and $u$-channels,
    respectively.  The last two diagrams in each of these lines involve the
    $\mathcal{O}(\lambda^2)$ vertex counterterm, while the final diagram is the
    $\mathcal{O}(\lambda^3)$ contribution to the vertex counterterm.}
  \label{fig:10.5}
\end{figure}

The value of the last diagram in \eqref{eq:10.51} is just a constant, which we
can freely adjust to absorb any divergent terms that are independent of the
external momenta.  Our goal, therefore, is to show that all momentum-dependent
divergent terms cancel among the remaining five diagrams.

The fourth and fifth diagrams in \eqref{eq:10.51} involve the one-loop vertex
counterterm, which we computed in Eq.~\eqref{eq:10.24}.  Let us briefly recall
that computation.  We defined $iV(p^2)$ as the fundamental loop integral,
\begin{equation}
  iV(p^2) = \frac{1}{2}\int \frac{d^dk}{(2\pi)^d}\ \frac{1}{k^2-m^2+i\epsilon}\
    \frac{1}{(k+p)^2-m^2+i\epsilon} .
  \label{eq:10.52}
\end{equation}
The counterterm, according to the renormalization condition \eqref{eq:10.19},
had to cancel the three one-loop diagrams (one for each channel) at threshold
($s = 4m^2$, $t = u = 0$); thus we found
\begin{equation}
  -i\delta_\lambda = (-i\lambda)^2\bigl[ -iV(4m^2) - 2iV(0) \bigr] .
\end{equation}
For our present purposes it will be convenient to separate the two terms of
this expression.  Let us therefore define
\begin{equation}
  V_s = (-i\lambda)^2\cdot(-iV(4m^2));\qquad
  V_u = (-i\lambda)^2\cdot(-2iV(0)) .
\end{equation}
We can now divide the first five diagrams in \eqref{eq:10.51} into three
groups:
\begin{equation}
  \text{Group I}:\quad \text{(two } s\text{-channel bubbles)};\qquad
  \text{Group II}:\quad \text{(one bubble + one counterterm)};\qquad
  \text{Group III}:\quad \text{(crossed version of Group II)} .
\end{equation}
We will find that all divergent terms that depend on momentum cancel
separately within each group.  Since Groups II and III are related by a simple
interchange of initial and final momenta, it suffices to demonstrate this
cancellation for Groups I and II.

Group I is actually quite easy, since each diagram factors into a product of
objects we have already computed.  Referring to Eq.~\eqref{eq:10.52}, we have
\begin{equation}
  \text{(bubble)}^2 = (-i\lambda)^3\bigl[ iV(p^2) \bigr]^2;
  \qquad \text{(bubble $\times$ counterterm)} = (-i\lambda)^3\
    \bigl[ -2iV(p^2)\,iV(4m^2) \bigr] .
\end{equation}
The sum of all three diagrams is therefore
\begin{equation}
  (-i\lambda)^3\Bigl( \bigl[ iV(p^2) \bigr]^2 - 2iV(p^2)\,iV(4m^2)
    \Bigr)
  = (-i\lambda)^3\Bigl( -\bigl[ V(p^2) - V(4m^2) \bigr]^2
    + \bigl[ V(4m^2) \bigr]^2 \Bigr) .
  \label{eq:10.53}
\end{equation}
But the difference $V(p^2) - V(4m^2)$ is finite, as was required for the
cancellation of divergences in the one-loop calculation:
\begin{equation}
  V(p^2) - V(4m^2) = \frac{1}{2}\int_0^1 dx\ \log\Bigl(
    \frac{m^2 - x(1-x)p^2}{m^2 - x(1-x)4m^2} \Bigr) .
\end{equation}
The only remaining divergence is in the term $[V(4m^2)]^2$, which is
independent of momentum and can therefore be absorbed into the second-order
counterterm in \eqref{eq:10.51}.

Two general properties of result \eqref{eq:10.53} are worth noting.  First,
the divergent piece (and hence the $\mathcal{O}(\lambda^3)$ vertex
counterterm) is proportional to
\begin{equation}
  \bigl[ V(4m^2) \bigr]^2 \propto \bigl[ \Gamma(2-d/2) \bigr]^2
  \qquad \text{for } d = 4-\epsilon .
\end{equation}
This is a \emph{double pole}, in contrast to the simple pole we found for the
one-loop counterterm.  Higher-loop diagrams will similarly have higher-order
poles, but in all cases the divergent terms are momentum-independent
constants.  Second, consider the large-momentum limit,
\begin{equation}
  V(p^2) - V(4m^2) \xrightarrow{p^2\to\infty}
    \frac{1}{2}\log^2\Bigl( \frac{p^2}{m^2} \Bigr) .
\end{equation}
The two-loop vertex is proportional to $\log^2(p^2/m^2)$.  A diagram of this
structure with $n$ loops will have the form
\begin{equation}
  \lambda^{n+1}\ \frac{1}{n!}\ \Bigl[ \log\Bigl( \frac{p^2}{m^2} \Bigr)
    \Bigr]^n .
\end{equation}
This asymptotic behavior is actually a generic property of multiloop diagrams,
which we will explore in more detail in Chapter~12.

Now consider the more difficult diagram, from Group II:
\begin{equation}
  \text{(bubble with vertex correction inserted)} .
  \label{eq:10.54}
\end{equation}
In evaluating this diagram, we will combine denominators in the manner that
makes it most straightforward to extract the divergent terms, at the price of
complicating the evaluation of the finite parts.  Another approach to the
calculation of this diagram is discussed in Problem~10.4.

To begin the evaluation of \eqref{eq:10.54}, combine the pair of denominators
shown explicitly, and substitute expression \eqref{eq:10.52} for $V(p^2)$.
This gives the expression
\begin{equation}
  \frac{(-i\lambda)^3}{2(4\pi)^{d/2}}\ \Gamma(2-d/2)\ \int_0^1 dx\ \int
    \frac{d^dk}{(2\pi)^d}\ \frac{1}{k^2-m^2}\ \frac{1}{(k+p)^2-m^2}\
    \frac{1}{\bigl[ m^2 - x(1-x)(k+p)^2 \bigr]^{2-d/2}} .
  \label{eq:10.55}
\end{equation}
It is possible to combine this pair of denominators by using the identity
\begin{equation}
  \frac{1}{A^\alpha B^\beta} = \frac{\Gamma(\alpha+\beta)}
    {\Gamma(\alpha)\Gamma(\beta)}\ \int_0^1 dw\ \frac{w^{\alpha-1}(1-w)^{\beta-1}}
      {\bigl[ wA + (1-w)B \bigr]^{\alpha+\beta}} .
  \label{eq:10.56}
\end{equation}
This is a special case of the formula quoted in Section~6.3,
Eq.~\eqref{eq:6.42}.  To prove it, change variables in the integral:
$w = x/(x+y)$, and use the definition of the beta function.

Carrying out the algebra, one finds that the divergent pieces of the Group II
diagram contain double poles and nonlocal logarithms, which exactly cancel
against the corresponding terms from the crossed diagram.  The remaining terms
are all either finite, or divergent but independent of momentum.  This
completes the proof that the two-loop contribution to the four-point function
is finite.

The two features of the Group I diagrams appear here in Group II as well.
The divergent pieces of the two-loop diagrams contain double poles that do not
cancel, so we again find that the second-order vertex counterterm must contain
a double pole.  The finite pieces contain double logarithms, so we again find
that the two-loop amplitude behaves as $\lambda^3\log^2 p^2$ as
$p\to\infty$.

\problems

\textbf{10.1} \emph{One-loop structure of QED.}  In Section~10.1 we argued
from general principles that the photon one-point and three-point functions
vanish, while the four-point function is finite.
\begin{enumerate}[label=(\alph*)]
  \item Verify directly that the one-loop diagram contributing to the
    one-point function vanishes.  There are two Feynman diagrams contributing
    to the three-point function at one-loop order.  Show that these cancel.
    Show that the diagrams contributing to any $n$-point photon amplitude, for
    $n$ odd, cancel in pairs.
  \item The photon four-point amplitude is a sum of six diagrams.  Show
    explicitly that the potential logarithmic divergences of these diagrams
    cancel.
\end{enumerate}

\textbf{10.2} \emph{Renormalization of Yukawa theory.}  Consider the
pseudoscalar Yukawa Lagrangian,
\begin{equation}
  \mathcal{L} = \tfrac{1}{2}(\partial_\mu\phi)^2 - \tfrac{1}{2}m^2\phi^2
    + \bar{\psi}(i\slashed{\partial} - M)\psi - ig\bar{\psi}\gamma^5\psi\phi ,
\end{equation}
where $\phi$ is a real scalar field and $\psi$ is a Dirac fermion.  Notice
that this Lagrangian is invariant under the parity transformation
\begin{equation}
  \psi(t,\mathbf{x}) \to \gamma^0\psi(t,-\mathbf{x}),\qquad
  \phi(t,\mathbf{x}) \to -\phi(t,-\mathbf{x}),
\end{equation}
in which the field $\phi$ carries odd parity.
\begin{enumerate}[label=(\alph*)]
  \item Determine the superficially divergent amplitudes and work out the
    Feynman rules for renormalized perturbation theory for this Lagrangian.
    Include all necessary counterterm vertices.  Show that the theory contains
    a superficially divergent $4\phi$ amplitude.  This means that the theory
    cannot be renormalized unless one includes a scalar self-interaction,
    \begin{equation}
      -\frac{\lambda}{4!}\phi^4 ,
    \end{equation}
    and a counterterm of the same form.  It is of course possible to set the
    renormalized value of this coupling to zero, but that is not a natural
    choice, since the counterterm will still be nonzero.  Are any further
    interactions required?
  \item Compute the divergent part (the pole as $d\to 4$) of each
    counterterm, to the one-loop order of perturbation theory, implementing a
    sufficient set of renormalization conditions.  You need not worry about
    finite parts of the counterterms.  Since the divergent parts must have a
    fixed dependence on the external momenta, you can simplify this
    calculation by choosing the momenta in the simplest possible way.
\end{enumerate}

\textbf{10.3} \emph{Field-strength renormalization in $\phi^4$ theory.}  The
two-loop contribution to the propagator in $\phi^4$ theory involves the three
diagrams shown in Eq.~(10.31).  Compute the first of these diagrams in
the limit of zero mass for the scalar field, using dimensional regularization.
Show that, near $d = 4$, this diagram takes the form:
\begin{equation}
  \sim \frac{\lambda^2}{16\pi^4}\ \frac{1}{4\epsilon}\ p^2 ,
\end{equation}
with $\epsilon = 4-d$.  The coefficient in this equation involves a Feynman
parameter integral that can be evaluated by setting $d = 4$.  Verify that the
second diagram in Eq.~(10.31) vanishes near $d = 4$.  Thus the first
diagram should contain a pole only at $\epsilon = 0$, which can be canceled by
a field-strength renormalization counterterm.

\textbf{10.4} \emph{Asymptotic behavior of diagrams in $\phi^4$ theory.}
Compute the leading terms in the $S$-matrix element for boson-boson scattering
in $\phi^4$ theory in the limit $s\to\infty$, $t$ fixed.  Ignore all masses on
internal lines, and keep external masses nonzero only as infrared regulators
where these are needed.  Show that
\begin{equation}
  i\mathcal{M} \sim -i\lambda - \frac{i\lambda^2}{32\pi^2}\ \log\Bigl(
    \frac{-s}{m^2} \Bigr) + \cdots .
\end{equation}
Notice that ignoring the internal masses allows some pleasing simplifications
of the Feynman parameter integrals.
